%!TEX root = TheoretischePhysikC.tex
% Kern von content/c03-hilbertraum, c03.s4 bis c03.s7.1. Quellen: docs/plan/script-v0.3/sync-ch3.md

\section{Vorlesung: Dirac-Notation, Hilbertraum, Basis}

\subsection{Dirac-Notation}
Sei $V$ ein komplexer Vektorraum; denken Sie vorerst an $\mathbb{C}^n$ aus der letzten Vorlesung. In der Quantenmechanik wird die folgende Schreibweise universell verwendet: Ab sofort schreiben wir jeden Vektor $\vb{v}\in V$ als \emph{Ket} $\ket{v}$. Den Fettdruck lassen wir im Ket weg; Strich und spitze Klammer sagen schon, dass ein Vektor gemeint ist. Mehr bedeutet die Notation zunächst nicht, und warum wir sie einführen, wird bald klar. Die eigentliche Hürde am Anfang ist nicht die Physik, sondern die Notation.

\begin{beispiel}
	Sei $V = \mathbb{C}^2$, der Raum, mit dem wir die Zustände eines Qubits beschreiben werden. Aus den beiden Vektoren
	\begin{equation}
		\ket{e_0} = \begin{pmatrix} 1 \\ 0 \end{pmatrix}, \qquad
		\ket{e_1} = \begin{pmatrix} 0 \\ 1 \end{pmatrix}
	\end{equation}
	bilden wir durch Addieren und Multiplizieren mit Zahlen zwei neue,
	\begin{equation}
		\ket{e_\pm} = \tfrac{1}{\sqrt{2}}\ket{e_0} \pm \tfrac{1}{\sqrt{2}}\ket{e_1}
		= \frac{1}{\sqrt{2}} \begin{pmatrix} 1 \\ \pm 1 \end{pmatrix}.
	\end{equation}
	Die Vektoren $\ket{e_\pm}$ tauchen immer wieder auf, beim Drehimpuls, beim Wasserstoffatom und besonders bei Spinsystemen.
\end{beispiel}

Dabei steht $\ket{v + w}$ für $\ket{v} + \ket{w}$ und $\ket{zv}$ für $z\ket{v}$, wenn $z$ eine Zahl ist.

Ergänzt wird die Ket-Notation durch die \emph{Bra}-Vektoren. Mit dem Skalarprodukt $\vb{v}\adj\vb{w}$ ist der Bra $\bra{v}$ die konjugierte Zeile, die die Kreuz-Operation aus der Spalte $\ket{v}$ macht. Das Produkt von Bra und Ket heißt \emph{Bra-Ket}; es ist das Skalarprodukt aus der letzten Vorlesung:
\begin{equation}
	\bra{v} = (\ket{v})\adj = \begin{pmatrix} v_1^* & \cdots & v_n^* \end{pmatrix}, \qquad
	\braket{v}{w} = (\ket{v})\adj\ket{w} = \vb{v}\adj\vb{w}.
\end{equation}
Das ist keine neue Definition, sondern eine neue Notation für dieselben Objekte.

\begin{historical}
	Die Notation stammt von Paul Dirac. Englisch \emph{bracket} heißt Klammer: links steht der ,,bra``, rechts der ,,ket``, und das ,,c`` ist irgendwo verlorengegangen.
\end{historical}

Ihren Vorteil spielt die Notation mit einer Regel für Linearkombinationen aus. Für komplexe Zahlen $\alpha,\beta$ gilt
\begin{equation}
	(\alpha\ket{v} + \beta\ket{w})\adj = \alpha^*\bra{v} + \beta^*\bra{w},
\end{equation}
denn die Kreuz-Operation konjugiert jeden Eintrag. Damit rechnen wir Skalarprodukte aus, ohne die Vektoren als Spalten hinzuschreiben. In unendlich vielen Dimensionen passt ein Vektor nicht in eine endliche Spalte, und dort rechnen wir genau so.

\begin{beispiel}
	Seien $\ket{v} = \ket{e_0} + i\ket{e_1}$ und $\ket{w} = \ket{e_0} + \ket{e_1}$. Nach der Regel ist $\bra{v} = \bra{e_0} - i\bra{e_1}$; beachten Sie das Minuszeichen. Mit $\braket{e_0}{e_0} = \braket{e_1}{e_1} = 1$ und $\braket{e_0}{e_1} = \braket{e_1}{e_0} = 0$ ist
	\begin{equation}
		\begin{split}
			\braket{v}{w} &= (\bra{e_0} - i\bra{e_1})(\ket{e_0} + \ket{e_1}) \\
			&= \braket{e_0}{e_0} + \braket{e_0}{e_1} - i\braket{e_1}{e_0} - i\braket{e_1}{e_1} = 1 - i.
		\end{split}
	\end{equation}
	Mit Spalten gerechnet ebenso: $1\cdot 1 + (-i)\cdot 1 = 1 - i$.
\end{beispiel}

Seien wir ehrlich: In zwei Dimensionen ist der erste Weg aufwendiger. In höheren Dimensionen werden Sie aber ganz schnell eine Liebe für die Ket-Notation entdecken.

\subsection{Skalarprodukte}
Auf $\mathbb{C}^n$ gibt es neben $\vb{v}\adj\vb{w}$ andere Skalarprodukte, sogar unendlich viele; die folgende Definition erfasst sie alle. Bei Rechnungen mit Spalten verwenden wir weiterhin $\vb{v}\adj\vb{w}$. Ein Skalarprodukt brauchen wir, weil die Quantenmechanik probabilistisch ist: Es ist unser Werkzeug für Wahrscheinlichkeiten.

\begin{definition}
	Ein \emph{Skalarprodukt} auf einem komplexen Vektorraum $V$ ordnet je zwei Vektoren $\ket{\phi},\ket{\psi}\in V$ eine komplexe Zahl $\braket{\phi}{\psi}$ zu. Es ist additiv in jedem Argument, $\braket{\phi_1 + \phi_2}{\psi} = \braket{\phi_1}{\psi} + \braket{\phi_2}{\psi}$ und ebenso im zweiten, und für alle $z\in\mathbb{C}$ gilt
	\begin{equation}
		\braket{\phi}{z\psi} = z\braket{\phi}{\psi}, \qquad \braket{z\phi}{\psi} = z^*\braket{\phi}{\psi}.
	\end{equation}
	Dazu kommt die \emph{Positivität}: $\braket{\psi}{\psi}$ ist reell und $\geq 0$, und null nur für $\ket{\psi} = 0$.
\end{definition}

Das Skalarprodukt $\vb{v}\adj\vb{w}$ erfüllt alle drei Bedingungen. Der Stern im ersten Argument ist wichtig: Wäre das Skalarprodukt dort linear, so wäre $\braket{i\psi}{i\psi} = -\braket{\psi}{\psi}$, und beide Seiten könnten nur für $\braket{\psi}{\psi} = 0$ nichtnegativ sein, auch wenn $\ket{\psi} \neq 0$ ist.

Die Positivität erlaubt uns, die Länge eines Vektors zu definieren, die \emph{Norm}
\begin{equation}
	\norm{\psi} = \sqrt{\braket{\psi}{\psi}} \geq 0.
\end{equation}
Mit dem Skalarprodukt $\vb{v}\adj\vb{w}$ ist $\norm{v}^2 = \vb{v}\adj\vb{v} = \sum_j\abs{v_j}^2$.

Aus den drei Bedingungen folgen die konjugierte Symmetrie und die \emph{Schwarz'sche Ungleichung}, aus dieser die Dreiecksungleichung für endlich viele Vektoren:
\begin{equation}
	\begin{gathered}
		\braket{\phi}{\psi} = \braket{\psi}{\phi}^*, \qquad
		\abs{\braket{\phi}{\psi}} \leq \norm{\phi}\,\norm{\psi}, \\
		\norm{\sum_{j=1}^{N}\ket{\psi_j}} \leq \sum_{j=1}^{N}\norm{\psi_j}.
	\end{gathered}
\end{equation}
Beide Ungleichungen sind Arbeitsgeräte. Die Schwarz'sche steckt hinter der Heisenbergschen Unschärferelation (Vorlesung~9); man beweist beide mit demselben Schritt: Das Skalarprodukt eines geschickt gewählten Vektors mit sich selbst ist nie negativ. Mit der Dreiecksungleichung schätzen wir Größen ab.

Zwei Vektoren heißen \emph{orthogonal}, $\ket{\phi} \perp \ket{\psi}$, wenn $\braket{\phi}{\psi} = 0$ ist. Achtung: Das heißt \emph{nicht}, dass einer der Vektoren null ist; das erzwingt die Positivität nur für $\ket{\phi} = \ket{\psi}$. Ein System von Vektoren $\ket{\phi_j}$ mit $j$ aus einer Indexmenge $X$ heißt \emph{orthonormal}, wenn
\begin{equation}
	\braket{\phi_j}{\phi_k} = \delta_{jk} \quad \text{für alle } j,k\in X.
\end{equation}
Dabei ist $\delta_{jk}$ das Kronecker-Delta, $1$ für $j = k$ und $0$ sonst; es ist nicht die Diracsche Delta-Funktion, die später vorkommt.

\subsection{Hilbertraum}
Jetzt haben wir genug beisammen, um Hilberträume zu definieren.

\begin{definition}
	Ein \emph{Hilbertraum} $\Hil$ ist ein komplexer Vektorraum mit Skalarprodukt, der bezüglich der Norm $\norm{\psi}^2 = \braket{\psi}{\psi}$ vollständig ist. \emph{Vollständig} heißt: Kommen die Glieder einer Folge $\ket{\psi_1},\ket{\psi_2},\dots$ in $\Hil$ einander beliebig nahe, $\norm{\psi_m - \psi_k} \to 0$ für $m,k \to \infty$, so hat die Folge einen Grenzwert $\ket{\psi}\in\Hil$, also $\norm{\psi_m - \psi} \to 0$.
\end{definition}

Die Vollständigkeit ist wichtig, aber subtil; sie hat mit Konvergenz in unendlichdimensionalen Räumen zu tun. In $\mathbb{C}^n$ mit $\vb{v}\adj\vb{w}$ ist sie automatisch erfüllt, und für Rechnungen dort ist ein Hilbertraum schlicht ein komplexer Vektorraum mit Skalarprodukt; für das Qubit ist das $\mathbb{C}^2$. In unendlich vielen Dimensionen gilt diese Abkürzung nicht: Dort braucht man die Vollständigkeit, etwa für die Entwicklung nach einer Basis mit unendlich vielen Vektoren.

\subsection{Basen}
Basen brauchen wir aus einem praktischen Grund: Die meisten Rechnungen der Quantenmechanik werden erst explizit, wenn eine Basis gewählt ist. Wir betrachten hier nur endliche Indexmengen $X$. Ein orthonormales System $\{\ket{\phi_j}\}_{j\in X}$ in $\Hil$ heißt \emph{Basis} (genauer \emph{Orthonormalbasis}) von $\Hil$, wenn es \emph{vollständig} ist:
\begin{equation}
	\norm{\psi}^2 = \sum_{j\in X}\abs{\braket{\phi_j}{\psi}}^2 \quad \text{für alle } \ket{\psi}\in\Hil.
\end{equation}
Verwirrenderweise sagt man auch hier vollständig, meint aber etwas anderes als bei der Vollständigkeit des Hilbertraums.

Die Paare $\ket{e_0},\ket{e_1}$ und $\ket{e_+},\ket{e_-}$ sind Basen von $\mathbb{C}^2$. Orthonormal ist auch das zweite Paar: $\braket{e_+}{e_-} = \tfrac12(1 - 1) = 0$ und $\norm{e_\pm}^2 = \tfrac12(1 + 1) = 1$. Für $\ket{\psi} = \alpha\ket{e_0} + \beta\ket{e_1}$ ist $\norm{\psi}^2 = \abs{\alpha}^2 + \abs{\beta}^2$; die rechte Seite ergibt dasselbe, mit $\braket{e_0}{\psi} = \alpha$ und $\braket{e_1}{\psi} = \beta$ ebenso wie mit $\braket{e_\pm}{\psi} = (\alpha \pm \beta)/\sqrt{2}$, denn $\abs{\alpha + \beta}^2 + \abs{\alpha - \beta}^2 = 2(\abs{\alpha}^2 + \abs{\beta}^2)$. In $\mathbb{C}^n$ mit $\vb{v}\adj\vb{w}$ gibt es immer eine Basis, die Spalten mit einer $1$ im $j$-ten Eintrag und sonst lauter Nullen. Im unendlichdimensionalen Fall ist alles schwieriger: Dass es eine Basis gibt, muss man beweisen, und der Beweis ist schwer. Basen mit unendlich vielen Vektoren behandelt der Webkurs.

Die wichtigste Eigenschaft einer Basis ist, dass man jeden Vektor $\ket{\Psi}\in\Hil$ nach ihr entwickeln kann, mit den Koeffizienten $c_j = \braket{\phi_j}{\Psi}$:
\begin{equation}
	\eqbox{\ket{\Psi} = \sum_{j\in X}\braket{\phi_j}{\Psi}\,\ket{\phi_j}.}
\end{equation}
Denn für $\ket{\xi} = \ket{\Psi} - \sum_j c_j\ket{\phi_j}$ ergibt die Orthonormalität $\braket{\xi}{\xi} = \norm{\Psi}^2 - \sum_j\abs{c_j}^2$, nach der Vollständigkeit null, und nach der Positivität ist $\ket{\xi} = 0$. Diese Gleichung ist eine Form der sogenannten \emph{Vollständigkeitsrelation}: Jeder Vektor ist die Summe seiner Komponenten längs der Basisvektoren, und die Koeffizienten sind Skalarprodukte. Sie sieht harmlos aus, wird aber ständig angewandt.

% fig_vl03_vollst.tex — Vollständigkeitsrelation: ein Vektor, zerlegt in
% zwei verschiedenen Orthonormalbasen desselben Raumes.
%
% Follows the conventions of fig_vl11_kasten.tex (the reference plate).
%
%   * geometry comes from the expansion derived in the lecture,
%         |Psi> = sum_j <phi_j|Psi> |phi_j>,
%     evaluated for the two orthonormal bases the lecture has already
%     introduced in C^2,
%         {|e_0>, |e_1>}   and   |e_±> = (|e_0> ± |e_1>)/sqrt(2),
%     on the concrete unit vector
%         |Psi> = (2|e_0> + |e_1>)/sqrt(5)   (real, so it fits in a plane).
%     Every coordinate below is that expansion, not a sketch:
%         c_0 = 2/sqrt5,  c_1 = 1/sqrt5   ->  |c|^2 = 4/5, 1/5
%         c_± = (2 ± 1)/sqrt10            ->  |c|^2 = 9/10, 1/10
%     and the two component rectangles therefore share |Psi> as their
%     diagonal exactly, because sum_j c_j |phi_j> = |Psi> in both bases.
%
%   * colour encodes *which basis*, and nothing else: cgblue for
%     {|e_0>,|e_1>}, birrenorange for {|e_+>,|e_->}.  The state |Psi>
%     itself is neutral navy — it belongs to no basis, which is the whole
%     point.  The coordinate rays are thin neutral grey: reference lines,
%     not data.
%
%   * within one basis the two members are told apart by tint (full
%     colour / 35%), not by a second hue, so the bars survive greyscale
%     printing; the numbers are printed as well, so nothing depends on
%     colour at all.  In (a) the same redundancy is carried by the dash
%     pattern — the first basis closes its rectangle with dashes, the
%     second with dots — and every arrow carries its own label, so the
%     two decompositions stay separable with the colour thrown away.
%     The neutral rays are held at birrengray!60: at full strength they
%     print at the same grey value as birrenorange and swallowed the
%     orange component vectors lying on top of them.
%
% Panel (a) the two decompositions of the same |Psi>.  Two different
%           rectangles, one diagonal.
% Panel (b) the weights |c_j|^2 those rectangles imply.  4/5 : 1/5 against
%           9/10 : 1/10 — the split is basis-dependent, the total is not.
%           This is the content of the completeness axiom
%           ||psi||^2 = sum_j |<phi_j|psi>|^2, and the reason the numbers
%           can later be read as probabilities.

\begin{center}
\begin{tikzpicture}[
    projvec/.style   ={-{Stealth[length=5pt,width=3.5pt]}, line width=1.3pt},
    statevec/.style  ={-{Stealth[length=6pt,width=4pt]}, line width=1.6pt,
                       spacecadet},
    basisray/.style  ={-{Stealth[length=4pt,width=3pt]}, line width=0.5pt,
                       birrengray!60},
    projdash/.style  ={line width=0.6pt, densely dashed},
    projdot/.style   ={line width=0.7pt, densely dotted},
]

% ─────────────────────────────────────────────────────────────
%  (a)  |Psi> resolved in two orthonormal bases
% ─────────────────────────────────────────────────────────────
\begin{axis}[
    name=geom,
    tpcplot,
    scale only axis,
    width=5.4cm, height=5.4cm,
    axis lines=none, grid=none,
    xmin=-0.38, xmax=1.38,
    ymin=-0.64, ymax=1.12,     % same span as x: 45 degrees looks like 45
    clip=false,
    title={(a) \tpcifenglish{One vector, two bases}%
                            {Ein Vektor, zwei Basen}},
]

  % ---- coordinate rays of the two bases (reference lines, neutral) ----
  \draw[basisray] (axis cs:-0.18,0)      -- (axis cs:1.24,0);
  \draw[basisray] (axis cs:0,-0.18)      -- (axis cs:0,1.00);
  \draw[basisray] (axis cs:-0.127,-0.127) -- (axis cs:0.721,0.721);
  \draw[basisray] (axis cs:-0.127,0.127)  -- (axis cs:0.424,-0.424);

  \node[lbl, anchor=west]       at (axis cs:1.27,0)      {$|e_0\rangle$};
  \node[lbl, anchor=south]      at (axis cs:0,1.03)      {$|e_1\rangle$};
  \node[lbl, anchor=south west] at (axis cs:0.74,0.74)   {$|e_+\rangle$};
  \node[lbl, anchor=north west] at (axis cs:0.44,-0.44)  {$|e_-\rangle$};

  % ---- the two component rectangles: same diagonal, different sides ----
  % basis {e_0,e_1}:  feet at (2/sqrt5, 0) and (0, 1/sqrt5)
  \draw[projdash, cgblue] (axis cs:0.894427,0)
                       -- (axis cs:0.894427,0.447214);
  \draw[projdash, cgblue] (axis cs:0,0.447214)
                       -- (axis cs:0.894427,0.447214);
  % basis {e_+,e_-}:  feet at 3/sqrt10 along e_+ and 1/sqrt10 along e_-
  \draw[projdot, birrenorange] (axis cs:0.670820,0.670820)
                            -- (axis cs:0.894427,0.447214);
  \draw[projdot, birrenorange] (axis cs:0.223607,-0.223607)
                            -- (axis cs:0.894427,0.447214);

  % ---- the component vectors c_j |phi_j> ----
  \draw[projvec, cgblue]       (axis cs:0,0) -- (axis cs:0.894427,0);
  \draw[projvec, cgblue]       (axis cs:0,0) -- (axis cs:0,0.447214);
  \draw[projvec, birrenorange] (axis cs:0,0) -- (axis cs:0.670820,0.670820);
  \draw[projvec, birrenorange] (axis cs:0,0) -- (axis cs:0.223607,-0.223607);

  % ---- the state itself ----
  \draw[statevec] (axis cs:0,0) -- (axis cs:0.894427,0.447214);
  \node[lbl, text=spacecadet, font=\small, anchor=south west]
    at (axis cs:0.91,0.47) {$|\Psi\rangle$};

  % ---- component labels, each beside its own arrow ----
  \node[lbl, text=cgblue, font=\scriptsize, anchor=north]
    at (axis cs:0.72,-0.045) {$c_0|e_0\rangle$};
  \node[lbl, text=cgblue, font=\scriptsize, anchor=east]
    at (axis cs:-0.04,0.22)  {$c_1|e_1\rangle$};
  \node[lbl, text=birrenorange, font=\scriptsize, anchor=south east]
    at (axis cs:0.58,0.685)  {$c_+|e_+\rangle$};
  \node[lbl, text=birrenorange, font=\scriptsize, anchor=north east]
    at (axis cs:0.155,-0.255) {$c_-|e_-\rangle$};

  \node[lbl, text=birrengray, font=\scriptsize, anchor=north east]
    at (axis cs:1.38,1.12) {$c_j \equiv \langle \phi_j|\Psi\rangle$};
\end{axis}

% ─────────────────────────────────────────────────────────────
%  (b)  The weights those two decompositions carry
% ─────────────────────────────────────────────────────────────
\begin{axis}[
    tpcplot,
    scale only axis,
    at={(geom.north east)}, anchor=north west, xshift=22mm,
    width=5.0cm, height=4.2cm,
    axis lines=left, grid=none,
    xmin=0, xmax=1.58, ymin=0.42, ymax=2.55,
    xtick={0,0.5,1}, xticklabels={$0$,$1/2$,$1$},
    ytick={1,2},
    yticklabels={{$\{|e_+\rangle,|e_-\rangle\}$},{$\{|e_0\rangle,|e_1\rangle\}$}},
    ytick style={draw=none},
    xlabel={$|c_j|^2$},
    clip=false,
    title={(b) \tpcifenglish{Weights}{Gewichte}},
]

  % the common total, drawn first so it does not sit on the bar borders
  \draw[levellinegrey] (axis cs:1,0.50) -- (axis cs:1,2.46);

  % basis {e_0,e_1}:  4/5 and 1/5
  \fill[cgblue]     (axis cs:0,1.74)  rectangle (axis cs:0.8,2.26);
  \fill[cgblue!35]  (axis cs:0.8,1.74) rectangle (axis cs:1.0,2.26);
  \draw[spacecadet, line width=0.5pt]
                    (axis cs:0,1.74)  rectangle (axis cs:1.0,2.26);

  % basis {e_+,e_-}:  9/10 and 1/10
  \fill[birrenorange]    (axis cs:0,0.74)  rectangle (axis cs:0.9,1.26);
  \fill[birrenorange!35] (axis cs:0.9,0.74) rectangle (axis cs:1.0,1.26);
  \draw[spacecadet, line width=0.5pt]
                         (axis cs:0,0.74) rectangle (axis cs:1.0,1.26);

  \node[font=\footnotesize, text=white] at (axis cs:0.40,2.00)
    {$|c_0|^2=\tfrac45$};
  \node[lbl, text=cgblue, font=\footnotesize, anchor=west]
    at (axis cs:1.04,2.00) {$|c_1|^2=\tfrac15$};

  \node[font=\footnotesize, text=white] at (axis cs:0.45,1.00)
    {$|c_+|^2=\tfrac{9}{10}$};
  \node[lbl, text=birrenorange, font=\footnotesize, anchor=west]
    at (axis cs:1.04,1.00) {$|c_-|^2=\tfrac{1}{10}$};

  \node[lbl, text=birrengray, font=\scriptsize]
    at (axis cs:0.50,1.50)
    {$\textstyle\sum_j |c_j|^2 = \|\Psi\|^2 = 1$};
\end{axis}

\end{tikzpicture}
\end{center}


Die Abbildung rechnet das für $\ket{\Psi} = (2\ket{e_0} + \ket{e_1})/\sqrt{5}$ in beiden Basen vor. Die Koeffizienten sind hier reell; nur deshalb passt (a) in eine Ebene. Jede Zerlegung schließt sich zu einem eigenen Rechteck: verschiedene Seiten, dieselbe Diagonale. In (b) stehen die Gewichte $\abs{c_j}^2$. Ihre Aufteilung hängt von der Basis ab, ihre Summe nicht: Sie ist $\norm{\Psi}^2$, für einen Vektor der Norm eins also eins.

Wie bildet man in der Praxis Basen? Mit dem Orthonormierungsverfahren von Gram und Schmidt. Aus einer beliebigen Liste $\ket{\Phi_1},\dots,\ket{\Phi_N}$ bestimmt man nacheinander
\begin{equation}
	\ket{\Psi_n} = \ket{\Phi_n} - \sum_{k=1}^{n-1}\frac{\braket{\Psi_k}{\Phi_n}}{\norm{\Psi_k}^2}\,\ket{\Psi_k},
\end{equation}
wobei Summanden mit $\ket{\Psi_k} = 0$ ausgelassen werden, weil sonst durch null geteilt würde. Dann ist $\braket{\Psi_k}{\Psi_n} = 0$ für $k<n$, und die Vektoren $\ket{\Psi_n}/\norm{\Psi_n}$ mit $\ket{\Psi_n} \neq 0$ bilden immer ein orthonormales System. Gibt es außer dem Nullvektor keinen Vektor, der auf allen $\ket{\Phi_n}$ senkrecht steht, so ist dieses System eine Basis von $\Hil$.

\begin{beispiel}
	In $\mathbb{C}^2$ starten wir mit zwei ausgedachten Vektoren,
	\begin{equation}
		\ket{\Phi_1} = \begin{pmatrix} 2 \\ 0 \end{pmatrix}, \qquad
		\ket{\Phi_2} = \begin{pmatrix} 1 \\ 1 \end{pmatrix}.
	\end{equation}
	Im ersten Schritt läuft die Summe von $1$ bis $0$; es wird gar nicht summiert, und $\ket{\Psi_1} = \ket{\Phi_1}$. Im zweiten Schritt ist mit $\braket{\Psi_1}{\Phi_2} = 2$ und $\norm{\Psi_1}^2 = 4$
	\begin{equation}
		\ket{\Psi_2} = \ket{\Phi_2} - \frac{\braket{\Psi_1}{\Phi_2}}{\norm{\Psi_1}^2}\,\ket{\Psi_1}
		= \begin{pmatrix} 1 \\ 1 \end{pmatrix} - \frac{1}{2}\begin{pmatrix} 2 \\ 0 \end{pmatrix}
		= \begin{pmatrix} 0 \\ 1 \end{pmatrix}.
	\end{equation}
	Geteilt durch die Normen $\norm{\Psi_1} = 2$ und $\norm{\Psi_2} = 1$ wird aus der ausgedachten Liste die Basis $\ket{e_0},\ket{e_1}$. Die Bedingung ist erfüllt, denn aus $\braket{\Phi_1}{\psi} = 2\psi_1 = 0$ und $\braket{\Phi_2}{\psi} = \psi_1 + \psi_2 = 0$ folgt $\ket{\psi} = 0$.
\end{beispiel}
